\documentclass[12pt,twoside]{article}
\usepackage[dvipsnames]{xcolor}
\usepackage{tikz,graphicx,amsmath,amsfonts,amscd,amssymb,bm,cite,epsfig,epsf,url}
\usepackage[hang,flushmargin]{footmisc}
\usepackage[colorlinks=true,urlcolor=blue,citecolor=blue]{hyperref}
\usepackage{amsthm,multirow,wasysym,appendix}
\usepackage{array,subcaption} 
% \usepackage[small,bf]{caption}
\usepackage{bbm}
\usepackage{pgfplots}
\usetikzlibrary{spy}
\usepgfplotslibrary{external}
\usepgfplotslibrary{fillbetween}
\usetikzlibrary{arrows,automata}
\usepackage{thmtools}
\usepackage{blkarray} 
\usepackage{textcomp}
\usepackage[left=0.8in,right=1.0in,top=1.0in,bottom=1.0in]{geometry}

\input{macros}

\begin{document}

\begin{center}
{\large{\textbf{Homework 5}} } \vspace{0.2cm}\\
Due  March 3 at 11 pm
\\
\end{center}
\input{hwstatement.tex}\\

\begin{enumerate}


\item (Modified bootstrap Gaussian confidence interval)
\begin{enumerate}
\item Explain how to modify the bootstrap Gaussian confidence interval for the correlation coefficient using Fisher's transformation. 
\item 
Apply your method on the correlation between \textit{footlength} and \textit{stature} in \texttt{ANSUR\ II\ MALE\ Public.csv} to compute 0.95 confidence intervals from $n$ samples. Plot them along with non-modified 0.95 bootstrap Gaussian confidence intervals and also with the population correlation coefficient.  You can try different sample sizes $n \in \{25, 100 \}$.
\item Estimate and report the probability that the population correlation belongs to each type of confidence interval. 
\end{enumerate}

\item (Confidence intervals for the variance) 
\begin{enumerate}
\item Explain how to build a $1-\alpha$ bootstrap Gaussian confidence intervals for the population variance from $n$ data points.
\item Explain how to build a $1-\alpha$ bootstrap percentile confidence intervals for the population variance from $n$ data points.
\item Build 0.95 bootstrap Gaussian and percentile confidence intervals for the population variance of \textit{stature} in \texttt{ANSUR\ II\ MALE\ Public.csv} setting $n=100$. Plot the confidence intervals along with the true population variance. 
\item Estimate and report the probability that the population variance belongs to each type of confidence interval.

\end{enumerate}

\end{enumerate}
\end{document}
